
Results demonstrate that benchmark saturation is algorithmically detectable with measurable expert consensus. Score convergence alone is sufficient for saturation detection (3/3 benchmarks, Levene's p < 0.05), though dual-metric integration (convergence + velocity) remains an architectural enhancement for future work. Temporal precedence analysis (100\% saturations preceded shifts by 32--78 months) distinguishes internal benchmark exhaustion from external paradigm disruption---saturation arises from intrinsic architectural exploration plateaus, not extrinsic events.

Expert consensus validation (76--93\% agreement) confirms saturation is a community-observable phenomenon, not an algorithmic artifact. Low-confidence response dispersion (30\% standard deviation) validates confidence scoring as a reliability indicator, justifying the filtering strategy for expert survey ground truth.

The magnitude of temporal precedence---mean 48-month lead time---far exceeds the original 6-month threshold prediction. This 4-year predictive window provides substantial early warning for proactive benchmark rotation. Additionally, the 46-month gap between algorithmic detection (ImageNet August 2015) and expert consensus \citep{June2019} suggests automated detection can provide early warning before community-wide recognition.
